\thispagestyle{empty}
\section*{Kurzdarstellung}
\label{sec:kurzdarstellung}
In dieser Arbeit wird die Effektivität von Smartphone-Sensordaten als Basis für die Klassifizierung von Bodenbelägen für Fahrradfahrer untersucht. Zur Klassifizierung werden mehrere Machine Learning Modelle getestet und verglichen. Über mehrere Monate wurden mit einem Citybike und einem handelsüblichen Smartphone die benötigten Test- und Trainingsdaten für unterschiedliche Straßenbeläge aufgenommen. Dabei kam die Applikation Phyphox zum Einsatz, die direkten Zugriff auf die Sensoren des Smartphones zur Aufnahme ermöglicht. Die Daten vorher händisch in die Kategorien Asphalt, Kopfstein, und Feldweg unterteilt, die die häufigsten Kategorien von Fahrradweg in einer europäischen Stadt darstellen. Zusätzlich werden sie in automatisch in Obeflächen-Rauheitsklassen der Stufen A bis F durch einen modifizierten International Roughness Index (IRI) - Alghorithmus eingeteilt. Am Ende hat sich herausgestellt, dass zwei Kombinationsmodelle die besten Ergebnisse lieferten: Ein Modell bestehend aus einem in Reihe geschalteten Convolutional Neural Network (CNN) und einem Long-Short-Term-Memory-Network (LSTM), sowie ein Modell aus einer in Reihe geschaltetem CNN und einem Transformer-Network. Beide erreichten eine Accuracy von 0,91.